\section{Contoh}
\label{sec:contoh-bab13}

Dua contoh berikut memakai percabangan dan perulangan pada program
utuh. Contoh pertama mencari nilai maksimum dengan \asm{LOOP} dan
\asm{JAE}. Contoh kedua mengurutkan array dengan perulangan
bersarang dan \asm{JBE}.

\subsection{Contoh 1: Mencari Nilai Maksimum}

Program pada Kode~\ref{lst:loop-maksimum} memeriksa setiap elemen
array dan menyimpan nilai terbesar pada \reg{AX}.

\lstinputlisting[
    language=MASM,
    caption={Program \file{loop-maksimum.asm} mencari nilai
    maksimum array.},
    label={lst:loop-maksimum}
]{\CodeDir/loop-maksimum.asm}

Alur percabangannya:

\begin{itemize}
    \item \lstasm|LEA SI, deret| menunjuk elemen pertama.
    \item \lstasm|MOV CX, BANYAK - 1| menyiapkan cacah perbandingan.
    \item \lstasm|CMP AX, [SI]| membandingkan kandidat dengan elemen
          saat ini.
    \item \lstasm|JAE LANJUT| melompat bila kandidat lebih besar atau
          sama. Lompatan ini memakai tafsiran tidak bertanda.
    \item \lstasm|MOV AX, [SI]| mengganti kandidat bila elemen lebih
          besar.
    \item \lstasm|LOOP CARI| mengulang sampai semua elemen diperiksa.
\end{itemize}

Keluaran program tampak pada Kode~\ref{lst:keluaran-maksimum}.

\begin{lstlisting}[
    style=sessiongaya,
    caption={Keluaran program \file{loop-maksimum.asm}.},
    label={lst:keluaran-maksimum}
]
Nilai maksimum: 93
\end{lstlisting}

\subsection{Contoh 2: Pengurutan Gelembung}

Program pada Kode~\ref{lst:bubble-sort} mengurutkan array byte
secara menaik dengan algoritme pengurutan gelembung.

\lstinputlisting[
    language=MASM,
    caption={Program \file{bubble-sort.asm} mengurutkan array
    dengan perulangan bersarang.},
    label={lst:bubble-sort}
]{\CodeDir/bubble-sort.asm}

Bagian penting program:

\begin{itemize}
    \item Perulangan luar memakai \reg{CX} untuk menentukan banyak
          lintasan.
    \item Pencacah luar disimpan dengan \asm{PUSH CX} sebelum
          perulangan dalam memakai \reg{CX}.
    \item Perulangan dalam membandingkan pasangan elemen yang
          berdekatan dengan \asm{CMP} dan \asm{JBE}.
    \item Bila elemen kiri lebih besar dari elemen kanan, kedua elemen
          ditukar.
    \item \asm{POP CX} memulihkan pencacah luar sebelum
          \asm{LOOP LINTASAN}.
\end{itemize}

Keluaran program tampak pada Kode~\ref{lst:keluaran-bubble}.

\begin{lstlisting}[
    style=sessiongaya,
    caption={Keluaran program \file{bubble-sort.asm}.},
    label={lst:keluaran-bubble}
]
Hasil pengurutan: 1 2 3 4 6 7 8 9
\end{lstlisting}

\begin{infobox}
\vspace{1.1em}
Kedua contoh memakai prosedur pembantu dan \eng{stack} untuk
pencetakan serta penyimpanan pencacah. Mekanisme \asm{CALL},
\asm{RET}, \asm{PUSH}, dan \asm{POP} dibahas pada bab berikutnya.
Pada bab ini, perhatian diarahkan pada alur \asm{CMP}, \asm{Jcc},
dan \asm{LOOP}.
\end{infobox}
